% PAPER_PHOTON_ARC_CQG --- standalone article. Originally targeted Classical and
% Quantum Gravity (CQG-116545, desk-rejected 07jul26 without review); revised and
% retargeted to General Relativity and Gravitation (Springer). Revision (jul/2026,
% post desk-rejection, wen_complement precedent): title now leads with the positive
% contributions (operator + H^2 law + map); the Campbell-Mecke valency sketch of
% sec:cst is promoted to a formal Proposition with proof sketch (source:
% docs/campaigns/IMPOSSIBILIDADE_PARCIAL/RESULTADO.md sec 3.1); abstract and intro
% reframed to lead with what the paper delivers. No number or claim changed.
%
% This is the "photon arc" paper: not a failure, a map. The TEIC programme asked
% repeatedly whether a photon (a massless transverse gauge vector, Coulomb phase)
% emerges from a causal set carrying an internal field. The answer is a sequence
% of pre-registered measurements, each relocating the obstruction:
%   E4   the orientation Goldstone modes are internal SCALARS, not a gauge vector
%        (polarisation does not lock to k; permutation p=0.23) -> photon excluded
%        in the orientation sector (this is the GoldstonePRD companion result).
%   E5/E7 the gauge-connection (link U(1)) sector: engine validated, but the bare
%        connection inherits the causal-set non-locality; the confine/Coulomb
%        phase shows a lattice-like crossover near beta~1 (Coulomb not certified).
%   E6-3 an indefinite-signature (E^2-B^2) Benincasa--Dowker gauge operator is
%        CONSTRUCTED (new; no precedent in the CST literature): gauge-invariant
%        to 5.7e-16, validated against free Maxwell on a 4D lattice (omega=ck,
%        c=1.05), but H2 FAILS on the causet -- causal diamonds are 100% electric
%        (B-type fraction = 0.0000), the magnetic sector is empty.
%   E6b  height of the diamond does NOT supply the spacelike 2-cell (tail ~0.25%,
%        non-growing).
%   E6c  spatial CURVATURE does: frac_B rises monotonically, crosses 0.01 at
%        R_hat=2 (b^2/e^2 per cell 0.11->0.97). The spacelike 2-cell exists and
%        grows with curvature -- the first positive lever.
%   E6d  coupling the ordered ferromagnet to the gauge field does NOT amplify it
%        (the rise is disorder noise; order suppresses it -- Meissner analogue
%        fails) -> dead.
%   E6e  extrapolation: Delta frac_B ~ (1/R_hat)^1.73 ~ H^2 (R^2=0.997); an O(10%)
%        magnetic sector needs R_hat~0.5-0.6 -- sub-Planckian. PHYSICAL frontier:
%        the mechanism is real but Planck-scale; it does NOT explain the photons
%        of the observable, low-curvature universe.
%
% NO photon is claimed as a positive result anywhere. The contribution is the
% systematic map + the new indefinite BD-gauge operator + the precise, original
% statement of the remaining obstruction (a spacelike Lorentz-invariant 2-cell at
% LOW curvature).
%
% Self-contained. Compiles with pdflatex+bibtex (revtex4-2 + apsrev4-2.bst)
% against TEIC_CONSOLIDATED.bib (copied into this folder). At submission this
% converts to the Springer class; the revtex build here is for review/length.
%
% Reframe (jul/2026, PLANO_REENQUADRAMENTO): intro positions the paper as the
% delimiting half of the CST matter-sector line (cites MatterGravityPRD + SU3PRD);
% sec:cst gains the structural (Campbell-Mecke / boost non-compactness) reason for
% the non-locality measured per-sector. No result or claim changed.
\documentclass[aps,prd,reprint,nofootinbib,superscriptaddress,floatfix]{revtex4-2}
\usepackage{amsmath,amssymb}
\usepackage{graphicx}
\usepackage{bm}
\usepackage{booktabs}
\usepackage{hyperref}

\newcommand{\Rhat}{\hat R}
\newcommand{\fracB}{f_{\!B}}

\newcounter{proposition}
\newenvironment{proposition}[1]{\par\medskip\noindent\refstepcounter{proposition}%
\textbf{Proposition~\theproposition}\ (\emph{#1})\textbf{.}\ \itshape}{\par\medskip}
\newenvironment{proofsketch}{\par\noindent\emph{Proof sketch.}\ }{\hfill$\square$\par}

\begin{document}

\title{An indefinite-signature Benincasa--Dowker gauge operator for causal sets:\\
electric causal diamonds, an $H^2$ curvature law, and the obstruction map\\
for an emergent photon}

\author{Miqueias Alves Mendes}
\email{mikalvesm@gmail.com}
\affiliation{Independent Researcher, Ibiapina, Cear\'a, Brazil}
\date{\today}

\begin{abstract}
Causal set theory possesses a smeared, Lorentz-invariant d'Alembertian for \emph{scalar}
fields, but no analogous operator for a gauge field. We construct one: an
indefinite-signature ($E^2-B^2$) Benincasa--Dowker gauge operator, in which each
causal-diamond plaquette enters the quadratic form with a sign set by its area bivector
(timelike/electric versus spacelike/magnetic). The operator is gauge-invariant to
$5.7\times10^{-16}$ and reproduces free Maxwell on a $4$D lattice ($\omega=ck$ with
fitted $c=1.05$; $c$ is never input). Applied to the Poisson causal set it returns a
sharp structural verdict: causal diamonds are $100\%$ electric (magnetic-cell fraction
$\fracB=0.0000$ exactly, one-sided Wilson bound $6\times10^{-5}$, stable for
$N=200$--$2000$)---the magnetic sector a Lorentzian signature requires is structurally
empty. With this tool we map, in a pre-registered elimination, where and why a photon is
hard to emerge on the causal set. (i)~The Goldstone modes of an internal orientation
order are two degenerate internal \emph{scalars}, not a gauge vector: their polarisation
does not lock to the wavevector (permutation $p=0.23$). (ii)~The gauge-connection (link
$U(1)$) sector is validated and not permanently confining (lattice-like crossover near
$\beta\simeq1$), but its bare connection inherits the causal set's non-locality (Coulomb
not certified)---and we prove this non-locality is a theorem, not an accident: by a
Campbell--Mecke (Palm) argument, \emph{every} Poincar\'e-invariant connection rule on
the sprinkling, pairwise or not, has divergent expected valency, because the boost orbit
at fixed proper time is the non-compact hyperboloid. (iii)~Diamond \emph{height} does
not supply the missing spacelike cell (tail $\sim0.25\%$, non-growing). (iv)~Spatial
\emph{curvature} does: $\fracB$ rises monotonically and crosses $0.01$ at de Sitter
radius $\Rhat=2$, and the curvature-added fraction obeys a clean law,
$\Delta\fracB\propto(1/\Rhat)^{1.73}\approx H^2$ ($R^2=0.997$). An $O(10\%)$ magnetic
sector then requires $\Rhat\approx0.5$--$0.6$---\emph{sub-Planckian}: the photon
mechanism is real but Planck-scale, and does not explain the photons of the observable
low-curvature universe. (v)~Coupling the ordered ferromagnet to the gauge field does not
amplify it (the rise is disorder, not order). (vi)~Two further low-curvature
levers---future-cone cells and a genuine Bianchi-I anisotropy on the causal
\emph{order}---are exhausted as well (the former gauge-trivial and off the Hasse-link
substrate, the latter redistributing but not magnetising), closing the low-curvature
map. The remaining obstruction is stated exactly: a Lorentz-invariant spacelike $2$-cell
on the causal order at low curvature.
\end{abstract}

\maketitle

\section{Introduction}\label{sec:intro}

That gauge fields might emerge as collective excitations of a structured vacuum---rather
than be postulated---is a recurring theme, from emergent photons in quantum spin liquids
and string-net condensates to the gauge sector of discrete quantum-gravity models
\cite{SverdlovBombelli2009}. For any such claim two things must be shown, not merely the
right \emph{number} of low-energy modes but the right \emph{tensorial character}: a
photon is a transverse \emph{vector} $A_\mu$ with a gauge redundancy, a Coulomb phase,
and a light-cone dispersion $\omega=ck$. The mode count alone never decides it.

We pose this for the discrete substrate of causal set theory (CST), where spacetime is
fundamentally a locally finite partial order---the order being the causal relation---and
a Lorentzian manifold is an approximation to it \cite{BLMS1987}. The substrate's
defining virtue is that a Poisson sprinkling breaks no Lorentz frame on average
\cite{BombelliHensonSorkin2009}, and CST recovers spacetime dimension and a
scalar-curvature (d'Alembertian) operator from the order alone
\cite{Myrheim1978,BenincasaDowker2010,Surya2019}. We take these geometric results as
established and add nothing to them; we ask only whether a \emph{photon} can be made to
emerge on the same substrate.

The paper's contribution is threefold. First, a new \emph{operator}: an
indefinite-signature Benincasa--Dowker gauge operator---to our knowledge the first
BD-type operator for a gauge $1$-form, and the first gated against free Maxwell on a
lattice before any causet reading (\S\ref{sec:operator}). Second, a
\emph{proposition}: the non-locality that obstructs every bare connection on the
sprinkling is forced by the non-compactness of the Lorentz group
(Proposition~\ref{prop:valency}, \S\ref{sec:theorem})---the measured obstructions of
\S\ref{sec:link} are instances of a theorem, not accidents of a construction. Third, a
\emph{map}: a pre-registered chain of measurements that relocates the obstruction to an
emergent photon step by step, ending at an open problem stated precisely enough to be
attacked. No photon is claimed anywhere; the negative is reported as carefully as the
positives, in the tradition of informative no-go results. The paper is also the
delimiting half of a matter-sector programme on the same substrate: companion
measurements derive an ordered vacuum with relativistic scalar Goldstone modes
\cite{GoldstonePRD}, a spin-$\tfrac12$ topological baryon that sources its own
gravitational field \cite{MatterGravityPRD}, and an $SU(3)$ colour sector with
confinement and a meson octet \cite{SU3PRD}; the present map bounds those
emergent-matter claims by a measured negative---the photon does not emerge in any
sector tested---rather than by silence. The programme that
generated it \cite{TEICDoc1} is cited once as provenance and never relied upon; every
claim below is a stand-alone measurement under a guard (\S\ref{sec:protocol}) that
forbids inserting relativity into the data-generating code. We proceed by elimination
through the natural carriers of a photon, recording at each step what is measured and
the obstruction that relocates the question:
the orientation Goldstone sector, where the modes are scalars (\S\ref{sec:goldstone});
the gauge-connection (link) sector, where non-locality blocks the bare connection
(\S\ref{sec:link}); an indefinite-signature Benincasa--Dowker gauge operator, new to
CST, which works on a lattice but finds the causal diamonds $100\%$ electric
(\S\ref{sec:operator}); and three levers on the missing magnetic cell---height, which
fails; curvature, which works; and an ordering coupling, which fails
(\S\ref{sec:levers}). We then quantify the one positive lever and find the photon
mechanism real but Planck-scale (\S\ref{sec:frontier}); we test the two further
low-curvature levers---future-cone cells and a genuine Bianchi-I anisotropy on the causal
order---and find both exhausted, closing the low-curvature map (\S\ref{sec:further});
discuss what this means for the
CST literature and prove that the non-locality met throughout is structural
(\S\ref{sec:cst}--\S\ref{sec:theorem}), and delimit precisely what is and is not
established (\S\ref{sec:claim}).

\section{Substrate and protocol}\label{sec:protocol}

The substrate is a causal set \cite{BLMS1987}: a finite set of events with a transitive,
irreflexive, locally finite partial order $\prec$, realised by Poisson sprinkling into
$(3{+}1)$-dimensional Minkowski space (or, in \S\ref{sec:levers}, de Sitter), with
$x\prec y$ when $y$ lies in the causal future of $x$. The \emph{links} are the relations
of the transitive reduction; a \emph{causal diamond} (order interval) is the set
$\{z: x\prec z\prec y\}$, and the elementary plaquettes of the gauge sector are built
from such diamonds. These are standard CST constructions, used without modification.

\emph{Anti-circularity.} Because we ask whether relativistic gauge structure
\emph{emerges}, we exclude having inserted it. A guard tokenises every data-generating
module and fails if a Lorentz factor, a light-cone speed, or any closed relativistic
form appears outside an explicitly labelled, non-generating comparison block. The speed
of light is never input: where a dispersion $\omega=ck$ is reported it is fitted to a
measured zero crossing, and $c$ appears only in quarantined comparison blocks. Every
number below is produced by code that passes this guard.

\section{The orientation Goldstone sector is scalar}\label{sec:goldstone}

The first natural carrier is the Goldstone mode of a spontaneously ordered internal
field. Placing a unit orientation $\vec n\in S^2$ on each event and coupling causal
neighbours by the minimal $O(3)$ tension
$S=J\sum_{\langle ij\rangle}\Delta\tau_{ij}[1-\vec n_i\!\cdot\!\vec n_j]$, the field
orders spontaneously through a continuous transition, and the order is genuine by
finite-size scaling (the order parameter rises with system size, Binder cumulant
$U_4=2/3$ throughout) \cite{GoldstonePRD}. The ordering $O(3)\to O(2)$ leaves
\emph{two} gapless transverse modes---the number a photon has.

But the number is not the character. A photon's polarisation is a spatial index locked
to the wavevector ($\vec k\!\cdot\!\vec\epsilon=0$); the two transverse modes here carry
an \emph{internal} $S^2$ index with, a priori, no relation to space, because the global
internal $O(3)$ is decoupled from the spacetime causal order. Measuring the
$2\times2$ polarisation tensor over $160$ wavevectors and testing the
eigenvector--$\hat k$ correlation against a confound-free permutation null, the
polarisation does \emph{not} track the wavevector (circular correlation $0.10$,
permutation $p=0.23$); the two modes are degenerate ($3.7\%$) and spatially isotropic
\cite{GoldstonePRD}. They are two internal \emph{scalar} Goldstone bosons, not a gauge
vector. This is not an accident of the model but the generic outcome: a $0$-form
Goldstone of a \emph{global internal} symmetry carries an internal index unlocked from
space, so it is structurally a scalar (in the model's non-Abelian extension these are
precisely the pseudoscalar mesons \cite{GellMann1962}, not a photon). A gauge vector
requires a field with a spacetime index \emph{and} a gauge redundancy---which an
internal order parameter does not provide. The photon, if it exists here, must live in a
different sector. We take this E4 result as the established starting point (it is the
subject of the companion paper \cite{GoldstonePRD}) and build forward from it.

\section{The gauge-connection sector and its non-locality}\label{sec:link}

The natural carrier with a genuine spacetime index and a gauge redundancy is a
\emph{connection} on the causal links: a $U(1)$ link variable $\theta_{ij}$ in the
spirit of lattice gauge theory \cite{Wilson1974}, whose Wilson loops around causal
diamonds define the field strength. Such a gauge \emph{action} has a precedent in CST
\cite{Sverdlov2008,SverdlovBombelli2009}: holonomies attached to causally related events
and a Lagrangian built from parallel transport around diamonds. What has never been done
is to measure its phase or its dispersion, and that is where the obstruction lives.

We validated the link engine---gauge invariance to machine precision, a $4$D ordering
transition, gauge-invariant diamond plaquettes---and then measured the
confinement/Coulomb phase via Wilson loops. The result is a \emph{crossover}
confine~$\to$~Coulomb near $\beta\simeq1$ that mirrors the compact-$U(1)$ lattice anchor
(the log-slope of large loops crosses from the confining band to the Coulomb band near
$\beta\approx0.85$, tracking the matched lattice points); a permanently confining $U(1)$
is therefore \emph{disfavoured}. But the Coulomb phase is \emph{not certified}: the clean
discriminator (the Creutz ratio on $R\times T$ rectangles) cannot be built on the
non-local causal order---there is no controlled rectangular loop structure---and a patch
surrogate fails its own validation control. The root cause is the documented non-locality
of causal-set operators \cite{Sorkin2007,BenincasaDowker2010}: the mean graph degree
grows without bound (empirically $\propto L^{2.9}$), and the region
$\{$future$,\ \tau\le\tau_{\max}\}$ has infinite volume in Minkowski, so there is no
finite-volume Lorentz-invariant local neighbourhood (\S\ref{sec:theorem} proves this is
forced for every invariant rule, not particular to the covering relation). The bare
connection inherits this;
a smeared operator is needed, exactly as the scalar sector needed the smeared
Benincasa--Dowker d'Alembertian.

A third, independent manifestation of the same non-locality shows that it obstructs not
merely the \emph{measurement} of the gauge phase but the gauge-generating
\emph{mechanism} itself. We tested the Anderson--Higgs route to a massive gauge field
directly: a gauge-invariant charged $U(1)$ condensate---an $O(3)$ order parameter gauged
under a non-compact $U(1)$ with covariant hopping---built on the same substrate. On a
regular $4$D lattice the mechanism is correct, the standard Higgs result: the gauge mass
grows with the condensate ($m_A:0\to0.63$ as the coupling turns on) and vanishes in the
disordered phase. On the causet it fails. Local spontaneous symmetry breaking needs the
charged condensate to hold a coherent phase across an event's neighbourhood, but the
causal order supplies $\sim25$ Hasse neighbours per event (against $8$ on the local
lattice, and growing with $N$), whose incoherent gauge phases frustrate it: the order
escapes to the neutral, ungauged axis (the charged bond order collapses $0.91\to0.00$
while the neutral order rises $0.00\to0.91$ already at coupling $\lambda=0.1$), $U(1)$ is
never broken, no Goldstone is absorbed, and the gauge field stays massless ($m_A=0$); a
stiffer gauge coupling does not rescue it. The Wilson-loop phase, the Coulomb
discriminator, and now the charged condensate are three independent observables blocked
by one structural fact: the non-locality of the Poisson causal order ($\sim25$ Hasse
neighbours per event) defeats any construction that requires local phase coherence. This
is a structural frontier of the substrate, not a technical limitation---and it is the
reason the next step is a \emph{smeared} operator, which tames exactly this non-locality.

\section{An indefinite-signature Benincasa--Dowker gauge operator}\label{sec:operator}

\subsection{Why an indefinite signature, and why it is new}
The scalar Benincasa--Dowker (BD) operator tames the non-locality with
alternating-sign layer weights that annihilate the divergent constant mode and converge
to $\Box$ in the continuum \cite{BenincasaDowker2010,DowkerGlaser2013}; its continuum
limit is proved \cite{BelenchiaBenincasaDowker2016}, which is the only reason it is
trusted. A literature scan (performed before any code) confirms that the entire BD /
Dowker--Glaser line defines the operator on a \emph{scalar field only}: there is no
BD-type operator for a vector, a connection, a $1$-form, or any $p$-form. The gauge line
in CST \cite{Sverdlov2008,SverdlovBombelli2009} supplies an \emph{action}, not a smeared
operator and not a dispersion. Moreover, a naive positive-definite gauge action
$\tfrac12\sum_P F_P^2$ is \emph{Euclidean}: it measures a non-negative energy with a
minimum at $\omega\approx0$, not a propagation. A Lorentzian photon needs an
\emph{indefinite} action, $E^2-B^2$, where the electric and magnetic modes enter with
opposite signs and produce the light cone.

We therefore construct an indefinite-signature BD-type gauge operator: each causal-
diamond plaquette is classified by its area bivector as \emph{E-type} (timelike bivector,
$A^{0i}$-dominant) or \emph{B-type} (spacelike, $A^{ij}$-dominant), and the two enter the
quadratic form with opposite signs,
\begin{equation}
S \;=\; \sum_{P\in E}\!F_P^2 \;-\!\sum_{P\in B}\!F_P^2 ,
\label{eq:indef}
\end{equation}
the $1$-form analogue of the scalar operator's alternating layer weights. This object has
no precedent in CST and no continuum theorem of its own; its correctness is part of the
research, and we gate it accordingly.

\subsection{Gauge invariance and the lattice validation gate}
Two checks precede any causal-set reading. \emph{Gauge invariance} is automatic
($F=d\theta$, so $\delta S=0$ under $\theta\to\theta+d\chi$ since $d^2=0$) and measured
to a relative $5.7\times10^{-16}$---the indefinite signature does not break it.
\emph{Continuum recovery}: on a regular $4$D lattice, where the free-Maxwell answer is
known and the E/B split is clean ($9408$ E-type, $9408$ B-type plaquettes on an $8^4$
lattice), the symbol's zero crossing follows $\omega=ck$ with a fitted light-cone speed
$c=1.05$ (deviation $3\%$), $c$ never inserted. The operator represents Maxwell where
Maxwell is known.

\subsection{The causet result: diamonds are 100\% electric}
On the Poisson causet the operator fails, and the reason is sharp. The causal diamonds
are \emph{entirely} E-type: the measured B-type fraction is $\fracB=0.0000$ exactly
across nine sprinklings ($N\le626$), and the exactly-zero value is stable as the
network grows---scanning $N\simeq200$ to $2000$ leaves $\fracB=0.0000$ with a
one-sided Wilson $95\%$ upper bound falling to $6\times10^{-5}$ over $6\times10^4$
plaquettes (Table~\ref{tab:fBN}), confirming the electricity of the diamonds is
structural, not a finite-size effect. Every causal diamond contains the timelike
extension from its past tip to its future tip, so its area bivector is timelike and
$b^2<e^2$ always; the magnetic sector that the $-\sum_B F_P^2$ term needs is
\emph{structurally empty}.

\begin{table}[t]
\caption{The exactly-zero magnetic ($B$-type) fraction is stable across network size
$N$. Each row pools three Poisson sprinklings at the target $N$; $\fracB$ is the
$B$-type fraction over all $P$ height-$2$ causal-diamond plaquettes, and the last
column is the one-sided Wilson $95\%$ upper bound for zero $B$-type cells out of $P$.
The structural electricity of the causal diamonds is not a finite-size effect.}
\label{tab:fBN}
\begin{ruledtabular}
\begin{tabular}{rrcc}
$N$ & plaquettes $P$ & $\fracB$ & Wilson $95\%$ upper \\
\colrule
$200$  & $6\,218$  & $0.0000$ & $6.2\times10^{-4}$ \\
$500$  & $59\,901$ & $0.0000$ & $6.4\times10^{-5}$ \\
$1000$ & $60\,000$ & $0.0000$ & $6.4\times10^{-5}$ \\
$2000$ & $60\,000$ & $0.0000$ & $6.4\times10^{-5}$ \\
\end{tabular}
\end{ruledtabular}
\end{table} The symbol is then $\lambda\propto-e^2\propto-\omega^2$,
strictly negative, and never crosses zero: no light cone. This is not the Euclidean
failure of the naive action (positive-definite, minimum at $\omega\approx0$); it is the
opposite failure of a correctly indefinite operator with \emph{no $B^2$ term to
balance}. The frontier is thereby \emph{relocated}, not crossed: the obstacle is neither
the signature (resolved) nor the BD weight (the normalised, sharp, and raw weights fail
identically), but the absence of a \emph{spacelike $2$-cell} in the causal order. That is
the precise open construction the rest of the paper probes.

\subsection{Continuum limit: an open mathematical problem}\label{sec:contlimit}
We are explicit about the status of the operator itself. The scalar Benincasa--Dowker
operator has a \emph{proven} continuum limit: it converges to $\Box$ as the sprinkling
density $\rho\to\infty$
\cite{BenincasaDowker2010,DowkerGlaser2013,BelenchiaBenincasaDowker2016}. The
indefinite-signature gauge extension constructed here is validated only
\emph{numerically}---on a regular $4$D lattice where Maxwell is known (gauge invariance
to $5.7\times10^{-16}$, $\omega=ck$ with $c=1.05$, \S\ref{sec:operator})---and its
continuum limit has \emph{not} been proven analytically. We regard this as an open
mathematical problem: the lattice validation is necessary but not sufficient for a
claim of mathematical equivalence to Maxwell. Accordingly we read the causet verdict as
a \emph{structural} statement about the causal-diamond $2$-complex (the cells are
electric), not as a proven continuum result---a distinction we keep throughout.

\section{Three levers on the spacelike 2-cell}\label{sec:levers}

If the obstruction is the absence of a spacelike $2$-cell, three geometric levers might
supply one: making the diamond taller, curving space, or coupling the ordered internal
field. We test each with a pre-registered success/death criterion, reusing the same
E/B classifier verbatim so that the levers are the only variable.

\subsection{Height does not help}\label{sec:height}
Generalising the height-$2$ plaquette to a $2h$-gon (two ascending chains of length $h$
between the tips; $h=2$ reproduces the original, B-type $=0.0000$, as an anchor), we
sweep $h=2$--$6$ at $N\le2000$. The exact zero is specific to height $2$. Spacelike
B-type cells \emph{appear} for $h\ge3$, but as a tail $\sim0.25\%$ that does \emph{not}
grow ($h{=}3{:}\,0.0024$, $h{=}6{:}\,0.0010$, declining at best sampling) and whose per-
cell $b^2/e^2$ \emph{decreases} with height ($0.22\to0.08$). No well-sampled cell reaches
the $0.01$ threshold. The route ``taller diamond'' to the spacelike $2$-cell is closed:
the magnetic sector is of measure zero at every height, with no light cone.

\subsection{Curvature does}\label{sec:curvature}
The orthogonal lever is spatial \emph{curvature}. A de Sitter background dS$_4$ in flat
slicing is conformally flat, so the \emph{causal order} in conformal coordinates is
identical to Minkowski's: sprinkling a conformal box gives the \emph{same} causet
ensemble at every curvature, with curvature the only variable (no loss of statistics).
The E/B split is read from the $5$D hyperboloid embedding $-X_0^2+\sum X_i^2=R^2$, where
curvature enters only through the area bivector. A mandatory gate passes: at $R\to\infty$
the curved construction reproduces the flat height scan bit-for-bit
($h{=}2{:}\,0.0000$, $h{=}3{:}\,0.0024$).

Curvature \emph{supplies} the magnetic sector, monotonically---the opposite of height.
The per-cell $b^2/e^2$ rises $8.6\times$ (from $0.11$ at the Minkowski floor toward $0.97$
under strong curvature), and the B-type fraction rises at every curvature step. At
$\Rhat\equiv R_{\rm dS}/\ell=2$ (de Sitter radius $\approx1.7\,\ell$), the measured
$\fracB=0.0117$ with a Wilson-style lower bound $0.0105>0.01$ over $\sim2.8\times10^4$
plaquettes (N-stable, not a fluctuation): the pre-registered success threshold fires.
The spacelike $2$-cell \emph{exists and grows with curvature}. We calibrate this
honestly: the crossing is marginal ($1.2\%$) and only at extreme curvature; the
\emph{average} diamond is still electric ($b^2/e^2=0.97<1$), so the magnetic sector is a
$\sim1.2\%$ tail at the edge of the light cone, not $O(1)$. (Height $2$ is immune,
$\fracB=0.0000$ at every curvature, since it always contains the tip-to-tip timelike
extension.) The obstruction is thus \emph{flat-specific and curvature-removable}; what
remains is to make the fraction $O(1)$, not a tail.

\subsection{An ordering coupling does not amplify it}\label{sec:coupling}
Before pushing curvature further, we test whether the ordered orientation field of
\S\ref{sec:goldstone} can \emph{amplify} the $1.2\%$ signal, by coupling it to the gauge
field, $\Phi=A+\lambda(\vec n_i\times\vec n_j)\!\cdot\!\hat e_z$, on the curved substrate
($\Rhat=2$). The coupling is realised exactly inside the E/B classifier ($\lambda=0$
reproduces \S\ref{sec:curvature} bit-for-bit). It does not amplify. The B-type fraction
does rise with $\lambda$ (to $0.27$ at $\lambda=2$)---but through \emph{disorder}, not
order: at $\lambda=2$ the disordered phase ($0.272$) exceeds the ordered phases
($J_c{:}\,0.205$, $2J_c{:}\,0.099$), so \emph{more order gives less magnetic fraction},
and at small $\lambda$ order even suppresses it below the baseline. Any incoherent
internal texture inflates $b^2$ as spacelike noise; ordering removes it. The Meissner-
like analogue (a coherent condensate amplifying the magnetic sector) does \emph{not}
occur. This lever is dead: the magnetic fraction can come only from curvature, not from
ordering.

\section{The mechanism is real, but Planck-scale}\label{sec:frontier}

Curvature is the one positive lever, so we ask quantitatively how much is needed for an
$O(1)$ photon. Fitting the five finite-curvature points (well-conditioned, parameter
uncertainties $\le17\%$), the curvature-added magnetic fraction obeys
\begin{equation}
\Delta\fracB \;=\; \fracB-\fracB^{\rm flat}\;\propto\;(1/\Rhat)^{1.73}\;\approx\;H^2,
\qquad R^2=0.997,
\label{eq:H2}
\end{equation}
i.e.\ the added fraction scales as the curvature squared (a fixed quadratic gives
$R^2=0.984$, as good within errors). Extrapolating, an $O(10\%)$ magnetic sector requires
$\Rhat\approx0.5$--$0.6$ (range $0.13$--$0.75$); $O(50\%)$ needs $\Rhat\approx0.2$--$0.3$;
even $O(5\%)$ needs $\Rhat\approx0.8$--$0.9$. Every target lies at $\Rhat<1$: a de Sitter
radius below the discreteness scale. Identifying the discreteness scale with the Planck
scale ($\ell\equiv\ell_{\rm Planck}$, declared external), this is \emph{sub-Planckian}
curvature---at or below the floor where the continuum/sprinkling picture is valid. A
pessimistic but comparably good exponential fit \emph{saturates} at $\sim3\%$, making
$O(10\%)$ unreachable at \emph{any} curvature.

The physical reading is unambiguous and is the central honest conclusion of the paper.
The curvature mechanism is real and cleanly characterised ($\fracB\propto H^2$,
$R^2=0.997$), but it delivers an $O(1)$ magnetic sector only at Planck-scale curvature.
The observable universe has $\Rhat$ astronomically large, where Eq.~\eqref{eq:H2} gives
$\fracB\approx$ the flat tail $\approx0$. \emph{The BD-gauge curvature route therefore
does not explain the photons of the observable, low-curvature universe.} This is a
\emph{physical} frontier, not merely a technical one: the ingredient (extreme curvature)
is well-defined and the scaling law is clean, but the required value lies where the
framework's own continuum description fails. A full Monte-Carlo campaign at $\Rhat<1$ is
of marginal value and not pursued; the well-conditioned extrapolation already answers the
question.

\section{Two further levers at low curvature}\label{sec:further}

Curvature is the only lever above that opens the magnetic sector, and only at the Planck
scale. Two further constructions remain that might supply a spacelike cell \emph{at low
curvature}; we test both under the same pre-registered discipline---with a control null in
each---and both fail, so the low-curvature map closes.

\subsection{Future-cone cells: gauge-trivial and off the link substrate}\label{sec:futurecone}
Instead of the diamond anchored on a timelike pair, anchor a cell on a \emph{spacelike}
pair ($i,j$ incomparable) with apices in their common future. The raw magnetic fraction is
spectacular---$\fracB=0.836$ at low curvature, N-stable, some $70\times$ the curved-diamond
value---but a control exposes it: a null with \emph{random} apices on the same spacelike
base already gives $0.614$, the kinematic floor of anchoring on a spacelike diagonal, so
only $\sim0.22$ comes from the causal structure. The pre-registered gate is gauge: for
$S=\tfrac12\sum F^2$ with $F=d\theta$ the boundary-of-boundary vanishes
($\partial\partial=0$) and the action is gauge-invariant to $\Delta S/S\approx10^{-15}$---
but this is \emph{automatic} for any set of closed quadrilaterals, and the kinematic null
reproduces the same machine zero, so the ``pass'' carries no information. Decisively, the
cells do not live on the substrate of Hasse links whose gauge field the link sector
validated: only $30.8\%$ of their sides are covering relations at $N=800$, and that fraction
\emph{falls} with $N$ ($0.452\to0.308$), the cells migrating off the link substrate as the
common-future apices deepen. They are a complex of causal relations, not a gauge complex;
the $\fracB=0.84$ was spacelike-base geometry, not a photon.

\subsection{Bianchi-I anisotropy on the causal order does not open it}\label{sec:bianchi}
The remaining low-curvature lever is genuine anisotropy placed in the \emph{causal order}
itself---directional light cones $\mathrm{d}t^2>a^2\,\mathrm{d}x_1^2+\mathrm{d}x_2^2+
\mathrm{d}x_3^2$ ($a=1$--$8$)---with the E/B split read in \emph{flat} coordinates (no
stretched axis, the error that made a coordinate-embedding version a pure artifact). The
magnetic fraction does not open: $\fracB\approx0.002$--$0.003$ for every $a$, never reaching
$0.01$, with a gap over an isotropic-order null of $\le0.0008$ (negative at $a=2$). The
anisotropy is not a trivial null---it genuinely \emph{redistributes} magnetic content,
draining the planes containing the squeezed axis ($b_{\rm plane}\,0.279\to0.011$) entirely
into the unsqueezed plane ($0.268\to0.629$)---but the diamonds remain timelike-dominated
(electric), so the \emph{fraction} with $b^2>e^2$ does not rise.

Together with isotropic curvature (Planck-scale, \S\ref{sec:curvature}), these exhaust the
three low-curvature levers; the magnetic sector stays a sub-percent electric-edge tail at
the curvature of the observable universe.\footnote{This causal-set non-locality is the same
high mean-field connectivity that, in the synthesis paper, leaves the ordering transition
without a divergent correlation length, so that no dimensional-transmutation scale is
manufactured from the substrate either.}

\section{What this means for the CST literature}\label{sec:cst}

The map sharpens an open problem in causal set theory into a single, concrete
construction. The Benincasa--Dowker / Dowker--Glaser operator line carries a smeared,
Lorentz-invariant d'Alembertian on \emph{scalars}
\cite{BenincasaDowker2010,DowkerGlaser2013,BelenchiaBenincasaDowker2016}; the gauge line
\cite{Sverdlov2008,SverdlovBombelli2009} supplies an \emph{action} but no smeared
operator, no measured phase, and no dispersion. The indefinite-signature gauge operator
of \S\ref{sec:operator} is, to our knowledge, the first BD-type operator for a gauge
$1$-form, and the first to be gated against free Maxwell on a lattice before any causet
reading. Its outcome localises the obstruction precisely: the $2$-complex of causal
diamonds carries only timelike (electric) cells, and curvature is the first ingredient
that endows it with spacelike (magnetic) cells, at a rate $\propto H^2$. The remaining
open problem is therefore not the signature, nor a different BD weight, nor confinement
(the link sector is not permanently confined, \S\ref{sec:link}), but the construction of
a \emph{Lorentz-invariant spacelike $2$-cell on the causal order at low curvature}---a
$2$-cell whose area bivector is spacelike without a preferred frame to call events
``simultaneous.'' Height, isotropic curvature, and an ordering coupling are exhausted
(dead, Planck-scale, and dead, respectively), and the two natural low-curvature directions
beyond them---a genuinely different $2$-cell (the future-cone construction,
\S\ref{sec:futurecone}) and a genuine Bianchi-I anisotropy on the causal order
(\S\ref{sec:bianchi})---are now tested and also fail: the former is gauge-trivial and lives
off the Hasse-link substrate, the latter redistributes magnetic content without raising its
fraction. The one genuinely open direction that remains is an \emph{inhomogeneous} geometry,
or a spacelike $2$-cell construction not yet conceived.

\section{The non-locality is a theorem}\label{sec:theorem}

The non-locality that blocks the bare connection (\S\ref{sec:link}) is, moreover, not an
accident of this construction. It is forced, for \emph{every} invariant connection rule
on the sprinkling, by the non-compactness of the Lorentz group. We state this precisely,
because it upgrades the measured obstructions of \S\ref{sec:link} from properties of one
construction to a structural fact about the substrate.

\begin{proposition}{invariant connection rules have divergent valency}\label{prop:valency}
Let $\Phi$ be a Poisson sprinkling of density $\rho$ in $d$-dimensional Minkowski space,
and let a connection rule join causally related events $i\prec j$ with a probability
that is a Poincar\'e-invariant measurable functional of $\Phi$ (pairwise, or depending
on any invariantly defined neighbourhood configuration). Then the expected valency of a
typical event is
\begin{equation}
\langle z\rangle \;=\; \rho\,\mathrm{Vol}\!\left(H^{d-1}\right)
\int_0^\infty \Delta\tau^{\,d-1}\, q(\Delta\tau)\,\mathrm{d}\Delta\tau ,
\label{eq:valency}
\end{equation}
where $q(\Delta\tau)$ is the connection probability marginalised over the rest of the
process (a function of the proper time $\Delta\tau$ alone, by invariance) and $H^{d-1}$
is the unit hyperboloid, the Lorentz orbit at fixed $\Delta\tau$. Since
$\mathrm{Vol}(H^{d-1})=\infty$, $\langle z\rangle=\infty$ unless $q\equiv0$ almost
everywhere (the trivially sparse regime, with no connections).
\end{proposition}

\begin{proofsketch}
By the Campbell--Mecke (Palm) formula for the Poisson process,
$\langle z\rangle=\rho\int_{J^+}h(i,x)\,\mathrm{d}^dx$, where
$h(i,x)=\mathbb{E}_\Phi\!\left[\mathbf 1(i,x\ \mathrm{connected})\right]$ is the
marginal connection probability. Both the sprinkling law and the rule are
Poincar\'e-invariant, so $h(g\cdot i,\,g\cdot x)=h(i,x)$ for every $g$ in the group;
hence $h$ depends only on the pair invariant $\Delta\tau$---including for rules that
depend on the whole neighbourhood configuration, whose $k$-point dependence is absorbed
into the marginal $h$. Decomposing the future-cone measure as
$\mathrm{d}^dx=\Delta\tau^{\,d-1}\,\mathrm{d}\Delta\tau\,\mathrm{d}\mu_{H^{d-1}}$ and
using that $h$ is constant on each hyperboloid (dependence on the position in $H^{d-1}$
would be frame dependence), the integral factorises and the infinite hyperbolic volume
exits as a prefactor.
\end{proofsketch}

Three remarks fix the scope. (a)~The covering relation itself is the best-known
instance: $q_{\rm cover}(\Delta\tau)=e^{-\rho c_d\Delta\tau^d}>0$, so links have
infinite expected valency in infinite Minkowski---the familiar non-locality of
causal-set links \cite{Sorkin2007,BenincasaDowker2010}; the content of
Proposition~\ref{prop:valency} is that this is not specific to links but holds for
every invariant rule, pairwise or not. (b)~In a finite box $\mathrm{Vol}(H^{d-1})$ is
regularised but grows with the system size (the accessible rapidity range grows), so
$\langle z\rangle$ diverges with $N$---exactly the $\propto L^{2.9}$ degree growth
measured in \S\ref{sec:link}. (c)~The nearest candidate exception, the
Alexandrov-interval abundance $N(i,j)$, carries no new spatial information: the interval
volume is fixed by $\Delta\tau$ itself, $V=c_d\Delta\tau^d$, so
$N(i,j)\sim\mathrm{Poisson}(\rho c_d\Delta\tau^d)$ and
$\mathbb{E}[N\,|\,\Delta\tau]=\rho c_d\Delta\tau^d$ is a monotone bijection---the same
variable in disguise (verified numerically, $\mathbb{E}[N]\propto\Delta\tau^3$ in
$2{+}1$D); every pairwise combinatorial quantity (longest-chain length, common-neighbour
counts) is likewise a function of $V\propto\Delta\tau^d$ and falls to the same
non-compactness. The smeared character of every workable causal-set
operator---including the gauge operator of \S\ref{sec:operator}---is the price of this
structural fact, not a technical choice.

\section{What is, and is not, established}\label{sec:claim}

\emph{Established by measurement.}
(a)~The orientation Goldstone modes are degenerate internal scalars, not a gauge vector
(polarisation unlocked from $\hat k$, $p=0.23$); a $0$-form Goldstone of a global
internal symmetry is structurally scalar (\S\ref{sec:goldstone}).
(b)~The link $U(1)$ sector is not permanently confining (lattice-like crossover near
$\beta\simeq1$); its bare connection inherits the causal-set non-locality, the clean
Coulomb discriminator is not constructible there, and an Anderson--Higgs charged
condensate is frustrated by the same non-locality ($U(1)$ unbroken, $m_A=0$, on a causet
where the regular-lattice Higgs mechanism works) (\S\ref{sec:link}).
(c)~An indefinite-signature BD gauge operator (new to CST) is gauge-invariant to
$5.7\times10^{-16}$ and reproduces free Maxwell on a $4$D lattice ($\omega=ck$,
$c=1.05$), but fails on the causet because causal diamonds are $100\%$ electric
($\fracB=0.0000$) (\S\ref{sec:operator}).
(d)~Diamond height does not supply the missing spacelike cell ($\sim0.25\%$,
non-growing) (\S\ref{sec:height}); curvature does ($\fracB$ monotonic, crossing $0.01$ at
$\Rhat=2$, $b^2/e^2$ per cell $0.11\to0.97$) (\S\ref{sec:curvature}); an ordering coupling
does not (the rise is disorder, order suppresses it) (\S\ref{sec:coupling}).
(e)~The curvature-added fraction scales as $\Delta\fracB\propto H^2$ ($R^2=0.997$), and an
$O(10\%)$ sector requires sub-Planckian curvature $\Rhat\approx0.5$--$0.6$
(\S\ref{sec:frontier}).
(f)~Two further low-curvature levers are exhausted: future-cone cells reach $\fracB=0.84$
but are gauge-trivial (the kinematic null reproduces the gauge invariance) and migrate off
the Hasse-link substrate ($30.8\%$ of sides are links at $N=800$, falling with $N$), and a
genuine Bianchi-I anisotropy on the causal order redistributes magnetic content between
spatial planes without raising $\fracB$ above $\sim0.003$ (\S\ref{sec:further}).

\emph{Established analytically.} Every Poincar\'e-invariant connection rule on the
Poisson sprinkling---pairwise or depending on any invariantly defined neighbourhood
configuration---has divergent expected valency, because the Lorentz orbit at fixed
proper time is the non-compact hyperboloid (Proposition~\ref{prop:valency}); the
non-locality met in \S\ref{sec:link} is structural to the substrate, not an artefact of
the constructions tested. The proposition's hypotheses are explicit (the invariant
Poisson measure, rules using only the order and invariant counts); it does not
constrain non-Poisson or dynamically grown measures.

\emph{Not established.} No photon is obtained: there is no causet regime, in any lever
tested, with an $O(1)$ transverse magnetic sector at the low curvature of the observable
universe, hence no emergent Maxwell field, no Coulomb phase certified, and no electric
charge or conserved current derived. The indefinite operator has no continuum theorem of
its own (it is validated only against the known lattice answer), so its causet verdict is
trusted as a \emph{structural} statement (the diamonds are electric) rather than a proven
continuum limit. The identification $\ell\equiv\ell_{\rm Planck}$ that turns ``$\Rhat<1$''
into ``sub-Planckian'' is an external assumption, declared. The Goldstone-scalar result
(a) is the companion paper's \cite{GoldstonePRD} and is used here as an established
input, not re-derived. This paper claims neither a photon nor its impossibility---it
claims a map, an operator, and a precisely located open construction.

\section{Conclusions}\label{sec:conclusions}

We have searched systematically for an emergent photon on a causal set and report the
search as an informative negative. The orientation Goldstone modes are internal scalars,
not a gauge vector, for a structural reason; the gauge-connection sector has the right
indices but inherits the substrate's non-locality---which independently frustrates a
charged Higgs condensate and leaves the gauge field massless; and an indefinite-signature
Benincasa--Dowker gauge operator---new to causal set theory, gauge-invariant and
correct against free Maxwell on a lattice---fails on the causet because the causal
diamonds are entirely electric. Of three geometric levers on the missing spacelike
magnetic cell, height fails, an ordering coupling fails, and only curvature works, at a
clean rate $\Delta\fracB\propto H^2$; but an $O(1)$ magnetic sector then demands
sub-Planckian curvature, so the mechanism, though real, does not produce the photons of
the low-curvature universe we inhabit. Two further low-curvature constructions---future-cone
cells and a genuine Bianchi-I anisotropy on the causal order---are then exhausted as well,
the first gauge-trivial and off the link substrate, the second redistributing magnetic
content without magnetising; the low-curvature map is closed. The contribution is not a photon but three portable
objects: a new gauge operator, validated where Maxwell is known; a proposition locating
the substrate's non-locality in the non-compactness of the Lorentz group, so that the
obstructions measured here are structural rather than constructional; and a chain of
pre-registered measurements that relocate the obstruction step by step to a single,
concrete open construction---a Lorentz-invariant spacelike $2$-cell on the causal order
at low curvature---on which an emergent photon in this framework now demonstrably
depends.

\begin{acknowledgments}
This work uses only the author's own simulations; the anti-circularity guard, the
validation gates, and the pre-registered success/death criteria of the E4--E6 and
Anderson--Higgs campaigns are part of the build.
\end{acknowledgments}

\bibliographystyle{apsrev4-2}
\bibliography{TEIC_CONSOLIDATED}

\end{document}
